\documentclass[10pt,a4paper]{article}

\usepackage[spanish]{babel}
\usepackage[margin=1.8cm]{geometry}
\usepackage{amsmath,amssymb}
\usepackage{multicol}
\usepackage{physics}
\usepackage{setspace}

\setlength{\parindent}{0pt}
\setlength{\parskip}{2pt}

\begin{document}

\begin{center}
{\Large \textbf{Identidades Trigonométricas y Complejas}}\\
{\small Chuleta para ondas, fasores y AC}
\end{center}

\hrule
\vspace{0.3cm}

\begin{multicols}{2}

%------------------------------------------------
\section*{Identidades fundamentales}

\[
\sin^2 x + \cos^2 x = 1
\]

\[
\tan x = \frac{\sin x}{\cos x}, \qquad
1+\tan^2 x = \frac{1}{\cos^2 x}
\]

\[
\cot x = \frac{\cos x}{\sin x}, \qquad
1+\cot^2 x = \frac{1}{\sin^2 x}
\]

%------------------------------------------------
\section*{Paridad y periodicidad}

\[
\sin(-x) = -\sin x, \qquad
\cos(-x) = \cos x
\]

\[
\sin(x+2\pi)=\sin x, \qquad
\cos(x+2\pi)=\cos x
\]

%------------------------------------------------
\section*{Suma y diferencia}

\[
\sin(a\pm b)=\sin a\cos b \pm \cos a\sin b
\]

\[
\cos(a\pm b)=\cos a\cos b \mp \sin a\sin b
\]

%------------------------------------------------
\section*{Doble, triple y medio ángulo}

\textbf{Doble ángulo}
\[
\sin(2x)=2\sin x\cos x
\]
\[
\cos(2x)=\cos^2 x-\sin^2 x
=2\cos^2 x-1=1-2\sin^2 x
\]

\textbf{Triple ángulo}
\[
\sin(3x)=3\sin x-4\sin^3 x
\]
\[
\cos(3x)=4\cos^3 x-3\cos x
\]

\textbf{Medio ángulo}
\[
\sin^2\frac{x}{2}=\frac{1-\cos x}{2}, \qquad
\cos^2\frac{x}{2}=\frac{1+\cos x}{2}
\]

%------------------------------------------------
\section*{Producto a suma}

\[
\sin a\sin b=\frac{1}{2}\Big[\cos(a-b)-\cos(a+b)\Big]
\]

\[
\cos a\cos b=\frac{1}{2}\Big[\cos(a-b)+\cos(a+b)\Big]
\]

\[
\sin a\cos b=\frac{1}{2}\Big[\sin(a+b)+\sin(a-b)\Big]
\]

%------------------------------------------------
\section*{Suma a producto}

\[
\sin a+\sin b=
2\sin\!\left(\frac{a+b}{2}\right)
\cos\!\left(\frac{a-b}{2}\right)
\]

\[
\cos a+\cos b=
2\cos\!\left(\frac{a+b}{2}\right)
\cos\!\left(\frac{a-b}{2}\right)
\]

%------------------------------------------------
\section*{Valores medios (ondas)}

\[
\langle \sin^2(\omega t)\rangle =
\langle \cos^2(\omega t)\rangle = \frac{1}{2}
\]

\[
\langle \sin(\omega t)\rangle =
\langle \cos(\omega t)\rangle = 0
\]

%------------------------------------------------
\section*{Identidades complejas}

\textbf{Fórmula de Euler}
\[
e^{ix}=\cos x+i\sin x
\]

\[
e^{-ix}=\cos x-i\sin x
\]

\[
\cos x=\frac{e^{ix}+e^{-ix}}{2}
\]

\[
\sin x=\frac{e^{ix}-e^{-ix}}{2i}
\]

%------------------------------------------------
\section*{Exponentes y fases}

\[
e^{i(a+b)}=e^{ia}e^{ib}
\]

\[
e^{i(\omega t+\phi)}=
\cos(\omega t+\phi)+i\sin(\omega t+\phi)
\]

\[
\Re\{e^{i\omega t}\}=\cos(\omega t), \qquad
\Im\{e^{i\omega t}\}=\sin(\omega t)
\]

%------------------------------------------------
\section*{Derivadas útiles}

\[
\frac{d}{dt}e^{i\omega t}=i\omega e^{i\omega t}
\]

\[
\frac{d^2}{dt^2}e^{i\omega t}=-\omega^2 e^{i\omega t}
\]

\end{multicols}

\vspace{0.2cm}
\hrule

\begin{center}
{\small Todo lo que necesitas para ondas, fasores y AC}
\end{center}

\end{document}
